\SourcePage{677}{680}
\section{The Rutherford formula}

Coulomb-field scattering is important in physical applications. It is also
noteworthy because, in this instance, the quantum-mechanical collision problem
admits a complete exact solution.

When a direction is distinguished---here, the direction of the incident
particle---the Schr\"odinger equation for a Coulomb field is conveniently
solved in the parabolic coordinates $\xi$, $\eta$, $\varphi$ (see \S~37). The
problem of scattering by a central field

\SourcePage{678}{681}
has axial symmetry, and the wave function $\psi$ is consequently independent
of the angle $\varphi$. Write a separated solution of the Schr\"odinger
equation (37.6) as
\begin{equation}
 \psi=f_1(\xi)f_2(\eta)
 \tag{135.1}
\end{equation}
(formula (37.7) with $m=0$). Separation of the variables then gives equations
(37.8) with $m=0$\footnote{Coulomb units are used throughout this section (see
p.~154).}:
\begin{align}
 \frac{d}{d\xi}\left(\xi\frac{df_1}{d\xi}\right)
 +\left[\frac{k^2}{4}\xi-\beta_1\right]f_1&=0, \notag\\
 \frac{d}{d\eta}\left(\eta\frac{df_2}{d\eta}\right)
 +\left[\frac{k^2}{4}\eta-\beta_2\right]f_2&=0, \notag\\
 \beta_1+\beta_2&=1.
 \tag{135.2}
\end{align}
The scattered particle naturally has positive energy; we have put $E=k^2/2$.
The signs in (135.2) refer to a repulsive field. An attractive field gives
exactly the same final scattering cross-section.

We require a solution of the Schr\"odinger equation which, for negative $z$
and large $r$, behaves as the plane wave
\[
 \psi\sim e^{ikz}\quad\text{for}\quad-\infty<z<0,\quad r\to\infty.
\]
This describes a particle incident in the positive $z$ direction. It will
follow below that a single separated solution (135.1), rather than a sum of
solutions having different $\beta_1$ and $\beta_2$, can meet this condition.

In parabolic coordinates the condition reads
\[
 \psi\sim\exp\left[\frac{ik}{2}(\xi-\eta)\right]
 \quad\text{for}\quad\eta\to\infty\quad\text{at every }\xi.
\]
It can hold only when
\begin{equation}
 f_1(\xi)=e^{ik\xi/2},
 \tag{135.3}
\end{equation}
while $f_2(\eta)$ obeys
\begin{equation}
 f_2(\eta)\sim e^{-ik\eta/2}\quad\text{for}\quad\eta\to\infty.
 \tag{135.4}
\end{equation}
Substitution of (135.3) into the first equation (135.2) confirms that this
function does solve it when $\beta_1=ik/2$. With
$\beta_2=1-\beta_1$, the second equation (135.2)

\SourcePage{679}{682}
therefore becomes
\[
 \frac{d}{d\eta}\left(\eta\frac{df_2}{d\eta}\right)
 +\left(\frac{k^2}{4}\eta-1+\frac{ik}{2}\right)f_2=0.
\]
Seek its solution in the form
\begin{equation}
 f_2(\eta)=e^{-ik\eta/2}w(\eta),
 \tag{135.5}
\end{equation}
where $w(\eta)$ tends to a constant limit as $\eta\to\infty$. The equation for
$w(\eta)$ is
\begin{equation}
 \eta w''+(1-ik\eta)w'-w=0.
 \tag{135.6}
\end{equation}
The change of variable $\eta_1=ik\eta$ reduces this to the confluent
hypergeometric equation with parameters $\alpha=-i/k$, $\gamma=1$. Among its
solutions we must choose the one which, after multiplication by $f_1(\xi)$,
contains only an outgoing (that is, scattered) spherical wave and no incoming
one. That solution is
\[
 w=\mathrm{const}\cdot F\left(-\frac{i}{k},1,ik\eta\right).
\]

Collecting these results, we obtain the following exact solution of the
Schr\"odinger equation for the scattering process:
\begin{equation}
 \psi=\exp\left(-\frac{\pi}{2k}\right)\Gamma\left(1+\frac{i}{k}\right)
 \exp\left[\frac{ik}{2}(\xi-\eta)\right]
 F\left(-\frac{i}{k},1,ik\eta\right).
 \tag{135.7}
\end{equation}
The normalization constant in $\psi$ has been selected so that the incident
plane wave has unit amplitude (as verified below).

To separate the incident and scattered waves in this function, we examine its
large-distance form. The first two terms in asymptotic expansion (d.14) of the
confluent hypergeometric function give, at large $\eta$,
\begin{align*}
 F\left(-\frac{i}{k},1,ik\eta\right)
 &\approx\frac{(-ik\eta)^{i/k}}{\Gamma(1+i/k)}
 \left(1+\frac{1}{ik^3\eta}\right)\\
 &\quad+\frac{(ik\eta)^{-i/k}e^{ik\eta}}{\Gamma(-i/k)ik\eta}\\
 &=\frac{e^{\pi/2k}}{\Gamma(1+i/k)}
 \left(1+\frac{1}{ik^3\eta}\right)
 \exp\left(\frac{i}{k}\ln k\eta\right)\\
 &\quad-\frac{(i/k)e^{\pi/2k}}{\Gamma(1-i/k)}
 \frac{e^{ik\eta}}{ik\eta}
 \exp\left(-\frac{i}{k}\ln k\eta\right).
\end{align*}
Insert this expression in (135.7) and change to spherical coordinates,
$\xi-\eta=2z$ and $\eta=r-z=r(1-\cos\theta)$. The resulting

\SourcePage{680}{683}
asymptotic form of the wave function is
\begin{align}
 \psi={}&\left[1+\frac{1}{ik^3r(1-\cos\theta)}\right]
 \exp\left\{ikz+\frac{i}{k}\ln[kr(1-\cos\theta)]\right\}\notag\\
 &+\frac{f(\theta)}{r}\exp\left\{ikr-\frac{i}{k}\ln(2kr)\right\},
 \tag{135.8}
\end{align}
where
\begin{equation}
 f(\theta)=-\frac{1}{2k^2\sin^2(\theta/2)}
 \frac{\Gamma(1+i/k)}{\Gamma(1-i/k)}
 \exp\left(-\frac{2i}{k}\ln\sin\frac{\theta}{2}\right).
 \tag{135.9}
\end{equation}

The first term of (135.8) is the incident wave. The Coulomb field falls off so
slowly that the incident plane wave remains distorted even far from the
center: the phase has a logarithmic contribution and the amplitude contains a
term of order $1/r$\footnote{This distortion is already intelligible from the
classical picture. Consider a family of classical hyperbolic Coulomb
trajectories with a common incident direction parallel to the $z$ axis. At
large distances from the scattering center ($z\to-\infty$), the equation of a
surface normal to these trajectories tends, as may readily be shown, not to
$z=\mathrm{const}$ but to
$z+k^{-2}\ln k(r-z)=\mathrm{const}$. This is precisely the constant-phase
surface of the incident wave in (135.8).}. The scattered spherical wave,
represented by the second term of (135.8), likewise has a logarithmic
distortion in its phase. These departures from the usual asymptotic form
(123.3) are nevertheless immaterial: their corrections to the flux density
vanish as $r\to\infty$.

Hence the scattering cross-section $d\sigma=|f(\theta)|^2d\Omega$ is
\[
 d\sigma=\frac{d\Omega}{4k^2\sin^4(\theta/2)},
\]
or, in ordinary units,
\begin{equation}
 d\sigma=\left(\frac{\alpha}{2mv^2}\right)^2
 \frac{d\Omega}{\sin^4(\theta/2)}.
 \tag{135.10}
\end{equation}
Here $v=k\hbar/m$ is the particle velocity. This is the familiar
\emph{Rutherford formula} of classical mechanics. Thus quantum and classical
mechanics predict the same result for scattering in a Coulomb field
(N.~Mott and W.~Gordon, 1928). The Born formula (126.12), as expected, also
leads to (135.10).

\SourcePage{681}{684}
For reference, we give the scattering amplitude (135.9) as an expansion in
spherical functions. Substitution in (124.5) of the phases from
(36.28) yields it\footnote{The quantity $\delta_l^{\mathrm C}$ in that formula
differs from the true Coulomb phase of the outgoing wave by a quantity common
to all $l$.}:
\begin{equation}
 \exp(2i\delta_l^{\mathrm C})=
 \frac{\Gamma(l+1+i/k)}{\Gamma(l+1-i/k)}.
 \tag{135.11}
\end{equation}
Consequently,
\begin{equation}
 f(\theta)=\frac{1}{2ik}\sum_l(2l+1)
 \frac{\Gamma(l+1+i/k)}{\Gamma(l+1-i/k)}P_l(\cos\theta).
 \tag{135.12}
\end{equation}

The signs in amplitude (135.9) correspond to repulsion in the Coulomb field.
For attraction, (135.9) must be replaced by its complex conjugate. In that
case $f(\theta)$ becomes infinite at the poles of $\Gamma(1-i/k)$, namely when
the argument of the gamma function is zero or a negative integer (here
$\Im k>0$, and $r\psi$ decays at infinity). The corresponding energies are
\[
 \frac{k^2}{2}=-\frac{1}{2n^2},\qquad n=1,2,3,\ldots,
\]
and coincide with the discrete energy levels of the attractive Coulomb field
(compare \S~128).
